\subsection{\texorpdfstring{第~\ref{sec:schwarzschild-solution}~节}{Schwarzschild 解}}\label{subsec:solutions-schwarzschild-solution}

\begin{description}[font=\normalfont,leftmargin=2em,labelwidth=1.5em]
\RGsolution{ex:equatorial-circular-geodesics-in-schwarzschild}{equatorial-circular-geodesics-in-schwarzschild}
\begin{enumerate}[label=(\alph*),ref=\alph*,leftmargin=*]
\item 设 $\{E_0,E_r,E_\theta,E_\varphi\}$ 为与标准正交余标架 $\{\omega^0,\omega^r,\omega^\theta,\omega^\varphi\}$ 对偶的标准正交标架, 即
\begin{align*}
  E_0&=\frac{1}{A}\frac{\partial}{\partial t},&
  E_r&=\frac{1}{B}\frac{\partial}{\partial r},\\
  E_\theta&=\frac{1}{r}\frac{\partial}{\partial\theta},&
  E_\varphi&=\frac{1}{r\sin\theta}\frac{\partial}{\partial\varphi}.
\end{align*}
由于
\[
  \dot c=\dot t\frac{\partial}{\partial t}+\dot\varphi\frac{\partial}{\partial\varphi}
        =\dot t A E_0+\dot\varphi r E_\varphi
\]
($\sin\theta=1$), 可得
\[
  \nabla_{\dot c}\dot c=0
  \quad\Longleftrightarrow\quad
  \frac{\mathrm{d}}{\mathrm{d}\tau}(\dot t A)E_0+\dot t A\nabla_{\dot c}E_0
    +\frac{\mathrm{d}}{\mathrm{d}\tau}(\dot\varphi r)E_\varphi
    +\dot\varphi r\nabla_{\dot c}E_\varphi=0.
\]
此外,
\[
  \nabla_{\dot c}E_0=\dot t A\nabla_{E_0}E_0+\dot\varphi r\nabla_{E_\varphi}E_0
\]
且
\[
  \nabla_{\dot c}E_\varphi=\dot t A\nabla_{E_0}E_\varphi+\dot\varphi r\nabla_{E_\varphi}E_\varphi.
\]
由于
\begin{align*}
  \nabla_{E_0}E_0&=\sum_{\mu=0}^3\omega_0^\mu(E_0)E_\mu
      =\omega_0^r(E_0)E_r=\frac{A'}{AB}E_r=A'E_r,\\
  \nabla_{E_0}E_\varphi&=\sum_{\mu=0}^3\omega_\varphi^\mu(E_0)E_\mu=0,\\
  \nabla_{E_\varphi}E_0&=\sum_{\mu=0}^3\omega_0^\mu(E_\varphi)E_\mu=0,\\
  \nabla_{E_\varphi}E_\varphi&=\sum_{\mu=0}^3\omega_\varphi^\mu(E_\varphi)E_\mu
      =\omega_\varphi^r(E_\varphi)E_r=-\frac{1}{Br}E_r,
\end{align*}
有
\[
  \nabla_{\dot c}E_0=\dot t AA'E_r
  \qquad\text{且}\qquad
  \nabla_{\dot c}E_\varphi=-\frac{\dot\varphi}{B}E_r,
\]
于是测地线方程成为
\[
  \ddot t A E_0+\left(\dot t^2A^2A'-\frac{\dot\varphi^2r}{B}\right)E_r
     +\ddot\varphi rE_\varphi=0,
\]
等价地,
\[
  \begin{cases}
    \ddot t=0\\
    \ddot\varphi=0\\
    \dot t^2A^2A'=\dfrac{\dot\varphi^2r}{B}
       \Longleftrightarrow AA'\dot t^2=\dot\varphi^2r
       \Longleftrightarrow \dfrac{m}{r^2}\dot t^2=\dot\varphi^2r.
  \end{cases}
\]
最后, 由于必须满足 $\langle\dot c,\dot c\rangle=-1$, 得
\begin{align*}
  -\left(1-\frac{2m}{r}\right)\dot t^2+r^2\dot\varphi^2=-1
    &\quad\Longleftrightarrow\quad
      \left(1-\frac{2m}{r}-\frac{m}{r}\right)\dot t^2=1\\
    &\quad\Longleftrightarrow\quad
      \left(1-\frac{3m}{r}\right)\dot t^2=1.
\end{align*}
\item 从上面的最后一个等式可见, 当 $r=3m$ 时, 有 $\langle\dot c,\dot c\rangle=0$, 因而得到一条赤道面上的圆形零模测地线. 若将一位静止观察者置于这条圆形光线上, 他会把它看成一条直线, 并看到自己无穷多个按相等间隔 (等于 $2\pi r$) 排列的像. 对应于这些像的光线在到达他的眼睛前绕了整数圈. 不同的像是这位观察者在不同时刻的像.
\item 设 $V=V^0E_0+V^rE_r+V^\theta E_\theta+V^\varphi E_\varphi$ 为一个绕点质量 $m$ 沿圆形轨道自由下落的自旋粒子的角动量向量场. 由于这个向量场与运动方向正交, 有
\begin{equation}\label{eq:schwarzschild-spin-orthogonal-to-orbit}
  -\dot t AV^0+\dot\varphi rV^\varphi=0.
\end{equation}
这个向量场沿运动轨道平行移动, 所以
\begin{align*}
  \nabla_{\dot c}V=0\quad\Longleftrightarrow\quad
   &\frac{\mathrm{d}V^0}{\mathrm{d}\tau}E_0+V^0\nabla_{\dot c}E_0
     +\frac{\mathrm{d}V^r}{\mathrm{d}\tau}E_r+V^r\nabla_{\dot c}E_r\\
   &+\frac{\mathrm{d}V^\theta}{\mathrm{d}\tau}E_\theta+V^\theta\nabla_{\dot c}E_\theta
     +\frac{\mathrm{d}V^\varphi}{\mathrm{d}\tau}E_\varphi+V^\varphi\nabla_{\dot c}E_\varphi=0.
\end{align*}
由于
\begin{align*}
  \nabla_{E_0}E_r&=\sum_{\mu=0}^3\omega_r^\mu(E_0)E_\mu=\omega_r^0(E_0)E_0=A'E_0,\\
  \nabla_{E_0}E_\theta&=\sum_{\mu=0}^3\omega_\theta^\mu(E_0)E_\mu=0,\\
  \nabla_{E_\varphi}E_r&=\sum_{\mu=0}^3\omega_r^\mu(E_\varphi)E_\mu
       =\omega_r^\varphi(E_\varphi)E_\varphi=\frac{1}{rB}E_\varphi,\\
  \nabla_{E_\varphi}E_\theta&=\sum_{\mu=0}^3\omega_\theta^\mu(E_\varphi)E_\mu=0,
\end{align*}
有
\begin{align*}
  \nabla_{\dot c}E_r&=\dot t A\nabla_{E_0}E_r+\dot\varphi r\nabla_{E_\varphi}E_r
      =\dot t AA'E_0+\frac{\dot\varphi}{B}E_\varphi,\\
  \nabla_{\dot c}E_\theta&=\dot t A\nabla_{E_0}E_\theta+\dot\varphi r\nabla_{E_\varphi}E_\theta=0,
\end{align*}
从而得
\begin{align}
  \frac{\mathrm{d}V^0}{\mathrm{d}\tau}&=-\dot t AA'V^r,\label{eq:schwarzschild-spin-time-component}\\
  \frac{\mathrm{d}V^r}{\mathrm{d}\tau}&=-\dot t AA'V^0+\frac{\dot\varphi}{B}V^\varphi,\notag\\
  \frac{\mathrm{d}V^\theta}{\mathrm{d}\tau}&=0,\label{eq:schwarzschild-spin-polar-component}\\
  \frac{\mathrm{d}V^\varphi}{\mathrm{d}\tau}&=-\frac{\dot\varphi}{B}V^r.\notag
\end{align}
将 \eqref{eq:schwarzschild-spin-orthogonal-to-orbit} 代入 \eqref{eq:schwarzschild-spin-time-component}, 得
\begin{equation}\label{eq:schwarzschild-spin-radial-component}
  \frac{\mathrm{d}V^r}{\mathrm{d}\tau}
   =\dot\varphi V^\varphi\left(\frac{1}{B}-A'r\right)
   =B\left(1-\frac{3m}{r}\right)\dot\varphi V^\varphi.
\end{equation}
所以, 由 \eqref{eq:schwarzschild-spin-polar-component} 和 \eqref{eq:schwarzschild-spin-radial-component}, 得
\begin{align*}
  \frac{\mathrm{d}V^r}{\mathrm{d}\varphi}
   &=\frac{\mathrm{d}V^r}{\mathrm{d}\tau}\frac{1}{\dot\varphi}
     =B\left(1-\frac{3m}{r}\right)V^\varphi,\\
  \frac{\mathrm{d}V^\varphi}{\mathrm{d}\varphi}
   &=\frac{\mathrm{d}V^\varphi}{\mathrm{d}\tau}\frac{1}{\dot\varphi}=-\frac{V^r}{B}.
\end{align*}
于是,
\[
  \frac{\mathrm{d}^2V^\varphi}{\mathrm{d}\varphi^2}
   =-\frac{1}{B}\frac{\mathrm{d}V^r}{\mathrm{d}\varphi}
   =-\left(1-\frac{3m}{r}\right)V^\varphi,
\]
从而
\[
  V^\varphi(\varphi)=\alpha\cos\left(\sqrt{1-\frac{3m}{r}}\,\varphi\right)
      +\beta\sin\left(\sqrt{1-\frac{3m}{r}}\,\varphi\right)
\]
且
\begin{align*}
  V^r(\varphi)&=-B\frac{\mathrm{d}V^\varphi}{\mathrm{d}\varphi}\\
    &=-B\sqrt{1-\frac{3m}{r}}
       \left(-\alpha\sin\left(\sqrt{1-\frac{3m}{r}}\,\varphi\right)
             +\beta\cos\left(\sqrt{1-\frac{3m}{r}}\,\varphi\right)\right).
\end{align*}
由于向量场 $V$ 起初沿径向, 有 $V^\varphi(0)=V^\theta(0)=V^0(0)=0$, 从而 $V^\theta\equiv0$, 并且
\begin{align*}
  V^\varphi(\varphi)&=\beta\sin\left(\sqrt{1-\frac{3m}{r}}\,\varphi\right),\\
  V^r(\varphi)&=-B\beta\sqrt{1-\frac{3m}{r}}
     \cos\left(\sqrt{1-\frac{3m}{r}}\,\varphi\right),\\
  V^0(\varphi)&=\frac{\dot\varphi r}{\dot t A}V^\varphi(\varphi)
     =\frac{\dot\varphi r}{\dot t A}\beta\sin\left(\sqrt{1-\frac{3m}{r}}\,\varphi\right).
\end{align*}
因此,
\begin{align*}
  V(0)&=-\beta B\sqrt{1-\frac{3m}{r}}\,X^r,\\
  V(2\pi)&=\frac{\dot\varphi r}{\dot t A}\beta
     \sin\left(2\pi\sqrt{1-\frac{3m}{r}}\right)X^0\\
    &\quad-\beta B\sqrt{1-\frac{3m}{r}}\cos\left(2\pi\sqrt{1-\frac{3m}{r}}\right)X^r
     +\beta\sin\left(2\pi\sqrt{1-\frac{3m}{r}}\right)X^\varphi,
\end{align*}
而由于
\[
  \langle V(0),V(2\pi)\rangle=\beta^2B^2\left(1-\frac{3m}{r}\right)
      \cos\left(2\pi\sqrt{1-\frac{3m}{r}}\right)
\]
且
\[
  |V(0)|^2=\beta^2B^2\left(1-\frac{3m}{r}\right)=|V(2\pi)|^2,
\]
得到
\[
  \cos(\delta)=\frac{\langle V(0),V(2\pi)\rangle}{|V(0)||V(2\pi)|}
      =\cos\left(2\pi\sqrt{1-\frac{3m}{r}}\right),
\]
其中 $\delta$ 是向量 $V(0)$ 和 $V(2\pi)$ 的夹角 (转过一周后). 选取 $\beta<0$, 使得 $V$ 起初与 $E_r$ 同向, 就容易验证 $\delta>0$, 所以
\[
  \delta=2\pi-2\pi\left(1-\frac{3m}{r}\right)^{\frac{1}{2}}.
\]
\end{enumerate}

\RGsolution{ex:radial-geodesics-in-schwarzschild}{radial-geodesics-in-schwarzschild}
\begin{enumerate}[label=(\alph*),ref=\alph*,leftmargin=*]
\item 由 Exercise~\ref{ex:schwarzschild-solution}(\ref{ex:schwarzschild-orbits-and-light-deflection}) 可见, 径向 ($\dot\theta=\dot\varphi=0$) 类时或零模测地线满足
\[
  \begin{cases}
    \dot r^2=E^2-\sigma\left(1-\dfrac{2m}{r}\right)\\[.6ex]
    \left(1-\dfrac{2m}{r}\right)\dot t=E
  \end{cases}
\]
其中 $E$ 为积分常数, 类时测地线取 $\sigma=1$, 零模测地线取 $\sigma=0$. 在这两种情形下, 满足 $r(0)=r_0$ 的测地线不能延拓到仿射参数
\[
  \lambda(2m)=\lambda(2m)-\lambda(r_0)
     =\int_{r_0}^{2m}\frac{\mathrm{d}\lambda}{\mathrm{d}r}\,\mathrm{d}r
     =\int_{r_0}^{2m}\frac{\mathrm{d}r}{\dot r}
\]
或
\[
  \lambda(0)=\lambda(0)-\lambda(r_0)
     =\int_{r_0}^0\frac{\mathrm{d}\lambda}{\mathrm{d}r}\,\mathrm{d}r
     =\int_{r_0}^0\frac{\mathrm{d}r}{\dot r}
\]
所给的值以外, 即不能超过
\[
  \pm\int_{r_0}^{2m}\frac{\sqrt r\,\mathrm{d}r}{\sqrt{(E^2-\sigma)r+2m\sigma}}
  \qquad\text{或}\qquad
  \pm\int_{r_0}^0\frac{\sqrt r\,\mathrm{d}r}{\sqrt{(E^2-\sigma)r+2m\sigma}}
\]
($\pm$ 为 $\dot r$ 的符号). 因为这两个积分显然都有限, 这些测地线必然不完备. 注意,
\[
  \int_{r_0}^{2m}\frac{\mathrm{d}t}{\mathrm{d}r}\,\mathrm{d}r
  =\int_{r_0}^{2m}\frac{\dot t}{\dot r}\,\mathrm{d}r
  =\int_{r_0}^{2m}\frac{Er^{\frac{3}{2}}\,\mathrm{d}r}{(r-2m)\sqrt{(E^2-\sigma)r+2m\sigma}}
\]
发散, 这意味着对一切满足 $E\ne0$ 的测地线, 当 $r\to2m$ 时, $t$ 发散 (只有在区域 $\{r<2m\}$ 中且 $t$ 不变的类时测地线才可能有 $E=0$).
\item 已知 $r$ 在仿射参数取某个有限值时趋向 $2m$. 由 Painlevé 时间坐标的定义可知, 它趋于
\begin{align*}
  \int_{r_0}^{2m}\frac{\mathrm{d}t'}{\mathrm{d}r}\,\mathrm{d}r
    &=\int_{r_0}^{2m}\frac{\mathrm{d}t}{\mathrm{d}r}\,\mathrm{d}r
     +\int_{r_0}^{2m}\frac{\sqrt{2mr}}{r-2m}\,\mathrm{d}r\\
    &=\pm\int_{r_0}^{2m}\frac{Er^{\frac{3}{2}}\,\mathrm{d}r}{(r-2m)\sqrt{(E^2-\sigma)r+2m\sigma}}
     +\int_{r_0}^{2m}\frac{\sqrt{2mr}}{r-2m}\,\mathrm{d}r\\
    &=\int_{r_0}^{2m}\frac{\sqrt{2mr}\sqrt{(E^2-\sigma)r+2m\sigma}\pm Er^{\frac{3}{2}}}
                         {(r-2m)\sqrt{(E^2-\sigma)r+2m\sigma}}\,\mathrm{d}r,
\end{align*}
其中 $r\to2m$ ($\pm$ 为 $\dot r$ 的符号). 考虑区域 $\{r>2m\}$ 内的测地线, 并选取 $E>0$, 对应于 $\dot t>0$. 此时, 当 $\dot r>0$ 时, Painlevé 时间坐标 $t'$ 发散; 当 $\dot r<0$ 时, 它收敛. 因此, 入射测地线可以越过 $r=2m$ 延拓, 而出射测地线不行. 换言之, 径向类时及零模测地线在过去渐近于 $r=2m$ (当仿射参数趋于某个有限值时, $t'$ 发散), 但在未来穿过该超曲面.
\item 由于
\[
  \frac{\partial}{\partial t'}
    =\frac{\partial t}{\partial t'}\frac{\partial}{\partial t}
     +\frac{\partial r}{\partial t'}\frac{\partial}{\partial r}
    =\frac{\partial}{\partial t},
\]
$\frac{\partial}{\partial t'}$ 仍然是 Killing 向量场. 因此, 径向曲线 $c:\mathbb{R}\to M$ 成为指向未来的类时测地线的方程可以写成
\[
  \begin{cases}
    g\left(\dfrac{\partial}{\partial t'},\dot c\right)=-E\\[.6ex]
    g(\dot c,\dot c)=-1
  \end{cases}
  \quad\Longleftrightarrow\quad
  \begin{cases}
    \dot t'-\sqrt{\dfrac{2m}{r}}\left(\dot r+\sqrt{\dfrac{2m}{r}}\dot t'\right)=E,\\[.6ex]
    \dot t'{}^2-\left(\dot r+\sqrt{\dfrac{2m}{r}}\dot t'\right)^2=1.
  \end{cases}
\]
所以, 只要
\[
  \dot t'=E=1,
\]
满足
\[
  \frac{\mathrm{d}r}{\mathrm{d}t'}=-\sqrt{\frac{2m}{r}}
     \quad\Longleftrightarrow\quad \dot r+\sqrt{\frac{2m}{r}}\dot t'=0
\]
的径向曲线就是指向未来的类时测地线. 特别地, 这意味着 Painlevé 时间坐标 $t'$ 与沿这些曲线的固有时一致.
\item 静止观察者接收到的光对应于出射的零模测地线. 如 (b) 所见, 这些测地线不能穿过视界, 因此像 Fig.~\ref{fig:stationary-observer-watching-fall-through-horizon} 中那样在视界处聚集. 于是它们都与表示下落粒子的曲线相交. 这意味着静止观察者将永远看到该粒子, 而粒子的运动越来越慢, 红移也越来越大.
\end{enumerate}
\end{description}
\RGfigure[.57\textwidth]{stationary-observer-watching-fall-through-horizon}{静止观察者看着粒子落入事件视界}
